% CHAPTER 5
\chapter{Results and Discussion}

\section{Experimental Setup}
[Recap the exact experimental conditions under which results were obtained: test partition size, hardware used for inference, batch size, and threshold parameters.]

\section{Training and Validation Curves}
[Present the accuracy, loss, or mAP curves across epochs. Comment clearly on model convergence, learning rate dynamics, and the presence or absence of overfitting/underfitting.]

\section{Performance Metrics on Test Set}
Table \ref{tab:test_metrics} presents the quantitative evaluation metrics computed on the unseen test dataset:

\begin{table}[htbp]
\centering
\caption{Performance metrics on the test dataset}
\label{tab:test_metrics}
\vspace{0.2cm}
\resizebox{0.95\columnwidth}{!}{
\begin{tabular}{|l|c|c|c|c|c|}
\hline
\textbf{Model / Configuration} & \textbf{Accuracy / mAP} & \textbf{Precision} & \textbf{Recall} & \textbf{F1-Score} & \textbf{Inference Speed} \\
\hline
[Baseline Model 1] & [Metric \%] & [Metric \%] & [Metric \%] & [Metric \%] & [ms / FPS] \\
\hline
[Baseline Model 2] & [Metric \%] & [Metric \%] & [Metric \%] & [Metric \%] & [ms / FPS] \\
\hline
\textbf{Proposed Model} & \textbf{[Metric \%]} & \textbf{[Metric \%]} & \textbf{[Metric \%]} & \textbf{[Metric \%]} & \textbf{[ms / FPS]} \\
\hline
\end{tabular}
}
\end{table}

\section{Confusion Matrix and Error Distribution}
[Provide the confusion matrix table or ROC curve analysis. Detail which specific classes are most frequently confused with each other, and analyze the underlying visual or mathematical reasons.]

\section{Comparison with Baseline Models}
[Contrast the proposed system against established standard baselines evaluated on identical test partitions (e.g., standard CNNs, classical ML baselines).]

\section{Comparison with Existing Works}
Table \ref{tab:lit_benchmarks} compares the proposed system against findings reported in recent published literature:

\begin{table}[htbp]
\centering
\caption{Comparison with published results from recent literature}
\label{tab:lit_benchmarks}
\vspace{0.2cm}
\resizebox{0.95\columnwidth}{!}{
\begin{tabular}{|l|l|l|l|}
\hline
\textbf{Study / Author} & \textbf{Core Methodology} & \textbf{Reported Metric} & \textbf{Key Limitation Addressed} \\
\hline
Author et al. \cite{1} & [Method 1] & [Metric 1] & [Limitation 1] \\
\hline
Author et al. \cite{2} & [Method 2] & [Metric 2] & [Limitation 2] \\
\hline
\textbf{Proposed System} & \textbf{[Proposed Method]} & \textbf{[Our Metric]} & \textbf{[Our Key Advantage]} \\
\hline
\end{tabular}
}
\end{table}

\section{Ablation Study}
[Detail the ablation experiments: systematically isolate or remove individual components (e.g., data augmentation, attention module, quantization, custom loss weighting) to quantify the individual performance contribution of each part.]

\section{Error Analysis}
[Show misclassified samples or failure edge cases. Honestly analyze the probable causes, such as severe occlusions, extreme lighting variations, sensor noise, or boundary ambiguities.]

\section{Discussion and Trade-Offs}
[Discuss what the results mean from an operational standpoint: highlight key strengths, acknowledge weak results honestly, and discuss practical trade-offs such as model accuracy versus inference latency or memory consumption.]
